\section{Expand sixteen words into sixty-four}
Each compression round uses one schedule word. There are 64 rounds, so the sixteen
words decoded from the block must be extended. The additional words combine earlier
schedule entries with two functions called small sigma:
\[
\begin{aligned}
\sigma_0(x)&=\operatorname{ROTR}_7(x)\mathbin{\oplus}
 \operatorname{ROTR}_{18}(x)\mathbin{\oplus}(x\mathbin{\gg}3),\\
\sigma_1(x)&=\operatorname{ROTR}_{17}(x)\mathbin{\oplus}
 \operatorname{ROTR}_{19}(x)\mathbin{\oplus}(x\mathbin{\gg}10).
\end{aligned}
\]
Here \(\oplus\) means bitwise XOR, \(\operatorname{ROTR}_n\) is a 32-bit right
rotation by \(n\), and \(\gg\) is a logical right shift. For indices 16 through 63,
compute
\[
W_t=W_{t-16}+\sigma_0(W_{t-15})+W_{t-7}+\sigma_1(W_{t-2})
\pmod{2^{32}}.
\]

Keep a pointer to the destination \(W_t\). Since a word occupies four bytes,
\(W_{t-16}\), \(W_{t-15}\), \(W_{t-7}\) and \(W_{t-2}\) are at byte offsets
\(-64\), \(-60\), \(-28\) and \(-8\) from that pointer. Those are byte offsets,
not schedule indices. The earlier entries have already been written before the loop
reads them. No initial contents of the schedule workspace should affect the result.

For the padded \texttt{abc} block, the first expanded values are
\texttt{61626380}, \texttt{000f0000}, \texttt{7da86405} and \texttt{600003c6}.
They are \(W_{16}\) through \(W_{19}\), in that order. If the first sixteen words
match but these do not, check the predecessor offsets and the distinction between a
rotation and a logical shift.

We retain all 64 words in this implementation so that the schedule is easy to inspect.
A circular sixteen-word schedule can reduce storage, but it changes which old entry is
overwritten at each iteration. That is a later experiment; the schedule interface in
this project deliberately exposes all 64 words for observation.
